\section{Lecture 1: Learning as Function Approximation}

\subsection{Motivation}

At its core, machine learning seeks to extract patterns from data in order to make predictions about unseen instances. This goal can be formulated in a unified and mathematically precise way by viewing learning as a problem of function approximation.

\medskip

\noindent
In many applications, we are given observations of the form:
\[
(x_i, y_i), \quad i = 1, \dots, n,
\]
where $x_i \in \mathcal{X}$ represents an input and $y_i \in \mathcal{Y}$ represents a corresponding output.

\medskip

\noindent
The central assumption is that there exists an (unknown) relationship between inputs and outputs:
\[
y = f(x),
\]
for some function $f : \mathcal{X} \to \mathcal{Y}$.

\medskip

\noindent
The task of learning is to approximate this unknown function using the observed data.

\medskip

\noindent
\textbf{Guiding idea.}  
Machine learning is the problem of constructing an approximation to an unknown function based on finite data.

\subsection{The Learning Problem}

We formalize the learning problem by introducing a parameterized family of functions:
\[
\{ f_\theta : \mathcal{X} \to \mathcal{Y} \mid \theta \in \Theta \},
\]
called a \emph{model class} or \emph{hypothesis space}.

\medskip

\noindent
Given data $\{(x_i, y_i)\}_{i=1}^n$, we seek parameters $\theta$ such that $f_\theta$ approximates the unknown function $f$.

\medskip

\noindent
To quantify approximation quality, we introduce a loss function:
\[
\ell : \mathcal{Y} \times \mathcal{Y} \to \mathbb{R}.
\]

\medskip

\noindent
The empirical learning problem becomes:
\[
\min_{\theta \in \Theta} \; \frac{1}{n} \sum_{i=1}^n \ell\big(f_\theta(x_i), y_i\big).
\]

\medskip

\noindent
\textbf{Interpretation.}
\begin{itemize}
    \item $f_\theta$ is our approximation to the unknown function $f$
    \item The loss measures discrepancy between predictions and observations
    \item Learning corresponds to selecting the best approximation within the model class
\end{itemize}

\subsection{Examples of Function Approximation}

\subsubsection{Linear Regression}

In linear regression, we assume:
\[
f_\theta(x) = w^\top x,
\]
where $\theta = w \in \mathbb{R}^d$.

\medskip

\noindent
Using squared loss:
\[
\ell(f(x), y) = (f(x) - y)^2,
\]
we obtain the objective:
\[
\min_w \sum_{i=1}^n (w^\top x_i - y_i)^2.
\]

\medskip

\noindent
This corresponds to approximating $f$ by a linear function.

\subsubsection{Classification}

In classification, outputs belong to discrete classes:
\[
y \in \{1, \dots, K\}.
\]

\medskip

\noindent
We may model:
\[
f_\theta(x) = \arg\max_k p_\theta(y = k \mid x).
\]

\medskip

\noindent
Here, the function approximation perspective remains valid: we approximate the mapping from inputs to labels.

\subsubsection{Nonlinear Models}

More expressive models introduce nonlinear transformations:
\[
f_\theta(x) = \sigma(Wx + b),
\]
or compositions of such transformations (as in neural networks).

\medskip

\noindent
These allow approximation of more complex functions.

\subsection{Approximation, Estimation, and Optimization}

The learning problem can be decomposed into three interacting components:

\medskip

\noindent
\textbf{1. Approximation.}  
The hypothesis space determines which functions can be represented.

\medskip

\noindent
\textbf{2. Estimation.}  
Given finite data, we estimate parameters $\theta$.

\medskip

\noindent
\textbf{3. Optimization.}  
We solve a numerical problem to find the best parameters.

\medskip

\noindent
\textbf{Insight.}  
Machine learning is not a single problem, but the interaction of approximation, estimation, and optimization.

\subsection{Geometry of Function Approximation}

It is often useful to interpret learning geometrically.

\medskip

\noindent
Let $\mathcal{H} = \{f_\theta\}$ denote the hypothesis space. Then learning can be viewed as finding:
\[
f_\theta \in \mathcal{H}
\]
that best matches observed data.

\medskip

\noindent
In linear regression, this corresponds to projecting the target vector $y$ onto the subspace spanned by input features.

\medskip

\noindent
\textbf{Geometric picture.}
\begin{itemize}
    \item Data defines a target point
    \item Model class defines a set (or manifold)
    \item Learning finds the closest point in that set
\end{itemize}

\subsection{Limitations of Finite Data}

Since we only observe finitely many samples, we cannot recover the true function $f$ exactly.

\medskip

\noindent
Two fundamental challenges arise:
\begin{itemize}
    \item \textbf{Underdetermination:} many functions may fit the data
    \item \textbf{Generalization:} performance on unseen data is not guaranteed
\end{itemize}

\medskip

\noindent
This motivates the need for:
\begin{itemize}
    \item appropriate model classes
    \item regularization
    \item probabilistic reasoning
\end{itemize}

\subsection{A Unifying Perspective}

Learning as function approximation provides a unifying framework for machine learning:

\begin{itemize}
    \item \textbf{Supervised learning:} approximate input-output mappings
    \item \textbf{Classification:} approximate decision functions
    \item \textbf{Regression:} approximate real-valued functions
    \item \textbf{Deep learning:} approximate complex compositions of functions
\end{itemize}

\medskip

\noindent
\textbf{Core abstraction.}  
All learning problems can be viewed as selecting a function from a hypothesis space to approximate an unknown target.

\medskip

\noindent
\textbf{Key insight.}  
The choice of function class determines what can be learned, while the data determines what is actually learned.

\subsection{Outlook}

In this lecture, we introduced the central paradigm of learning as function approximation.

\medskip

\noindent
However, this formulation raises several important questions:

\begin{itemize}
    \item How is the data generated?
    \item What does it mean for a model to generalize?
    \item How should we choose the loss function?
\end{itemize}

\medskip

\noindent
In the next lecture, we address these questions by adopting a probabilistic perspective, where data is modeled as random variables and learning is framed as statistical estimation.

\medskip

\noindent
\textbf{Guiding principle.}  
Understanding learning requires combining approximation theory with probabilistic reasoning.